%%%%%%%%%%  scan idx 34  |  folio 21  %%%%%%%%%%
\section*{CHAPITRE I}
\begin{center}
$\lambda$-Anneaux (Préliminaires formels)
\end{center}

\subsection*{\S 1. Définitions}

On appelle \emph{$\lambda$-anneau}\footnote{Dans ce séminaire, nous disons plutôt \emph{pré-$\lambda$-anneau} (V 2.1), réservant le nom de $\lambda$-anneau à ceux satisfaisant la condition supplémentaire de la page 3 (cf V 2.4).} un anneau commutatif $K$ avec unité, muni d'une famille d'applications
\begin{equation}\tag{1.1}
\lambda^i:K\longrightarrow K
\end{equation}
($i$ entier $\geq 0$) satisfaisant aux conditions suivantes~:
\begin{equation}\tag{1.2}
\lambda^0x=1,\qquad \lambda^1x=x,\qquad
\lambda^n(x+y)=\sum_{i=0}^n\lambda^ix\cdot\lambda^{n-i}y.
\end{equation}
Posons, pour tout $x\in K$
\begin{equation}\tag{1.3}
\lambda_t(x)=\sum_{n\geq 0}\lambda^n(x)t^n\in K[[t]] ;
\end{equation}
les relations précédentes expriment alors que $x\longmapsto\lambda_t(x)$ est un homomorphisme du groupe additif $K$ dans le groupe multiplicatif $1+K[[t]]^+$ des séries formelles dans $K$, d'augmentation $1$, relevant l'homomorphisme canonique $1+\sum_{i>0}x_it^i\longrightarrow x_1$.

Ainsi, sur l'anneau $\Z$ existe une unique $\lambda$-structure compatible avec sa structure d'anneau et telle que
